% Figure S2: the four split designs as host sets, names only. Caption: fig_splits.caption.md
\documentclass[tikz,border=1pt]{standalone}
% Shared preamble for the TikZ schematics (fig_*.tex). Load with % Shared preamble for the TikZ schematics (fig_*.tex). Load with % Shared preamble for the TikZ schematics (fig_*.tex). Load with \input{common} after
% \documentclass[tikz,border=1pt]{standalone}.
% Palette = src/dftgnn/viz/style.py COLORS (Okabe-Ito hues), extended with P and the v2 style.
\usepackage[T1]{fontenc}
\usepackage[scale=1.0]{tgheros}            % Helvetica-metric sans, Type 1, embeddable
\renewcommand{\familydefault}{\sfdefault}
\usepackage[eulergreek]{sansmath}
\sansmath
\usepackage{tikz}
\usetikzlibrary{positioning,arrows.meta,fit,backgrounds,calc,shapes.geometric,decorations.pathreplacing}

% ---- palette (hex as in style.py) ----
\definecolor{cB0}{HTML}{7F7F7F}            % grey   B0 physics floor
\definecolor{cRF}{HTML}{E69F00}            % orange RF-Kumagai
\definecolor{cS}{HTML}{0072B2}             % blue   S
\definecolor{cD}{HTML}{009E73}             % green  D
\definecolor{cP}{HTML}{CC79A7}             % reddish purple P (P1 and the staged P)
\definecolor{cRef}{HTML}{4D4D4D}           % reference grey (arrows, neutral outlines)
\definecolor{cO}{HTML}{D55E00}             % vermillion: the flagged vacancy O site
\colorlet{cNeutral}{black!45}

% ---- page geometry (docs/venue.md: MDPI \textwidth 138.6 mm) ----
\newlength{\figwidth}\setlength{\figwidth}{138.6mm}

% ---- styles ----
\tikzset{
  every picture/.style={line width=0.6pt, font=\fontsize{8}{9.6}\selectfont},
  >={Stealth[length=2.2mm,width=1.6mm]},
  flow/.style={->, draw=cRef, line width=0.6pt},
  inject/.style={->, draw=cD, line width=0.9pt},
  box/.style={draw=cNeutral, rounded corners=1.2pt, minimum height=6mm, inner xsep=2.2pt,
              inner ysep=1.6pt, align=center, fill=white},
  model/.style={box, draw=#1, fill=#1!14, line width=0.8pt},
  modelname/.style={font=\fontsize{9}{10}\selectfont\bfseries, text=#1},
  quantity/.style={box, draw=cNeutral, fill=black!4},
  v2/.style={draw=#1, dashed, dash pattern=on 2pt off 1.2pt, line width=0.8pt, rounded corners=1.2pt},
  panel/.style={font=\fontsize{9}{10}\selectfont\bfseries, anchor=north west, inner sep=0pt},
  atom/.style={circle, inner sep=0pt, minimum size=#1},
}
\newcommand{\Ef}{$E_\mathrm{f}$}
 after
% \documentclass[tikz,border=1pt]{standalone}.
% Palette = src/dftgnn/viz/style.py COLORS (Okabe-Ito hues), extended with P and the v2 style.
\usepackage[T1]{fontenc}
\usepackage[scale=1.0]{tgheros}            % Helvetica-metric sans, Type 1, embeddable
\renewcommand{\familydefault}{\sfdefault}
\usepackage[eulergreek]{sansmath}
\sansmath
\usepackage{tikz}
\usetikzlibrary{positioning,arrows.meta,fit,backgrounds,calc,shapes.geometric,decorations.pathreplacing}

% ---- palette (hex as in style.py) ----
\definecolor{cB0}{HTML}{7F7F7F}            % grey   B0 physics floor
\definecolor{cRF}{HTML}{E69F00}            % orange RF-Kumagai
\definecolor{cS}{HTML}{0072B2}             % blue   S
\definecolor{cD}{HTML}{009E73}             % green  D
\definecolor{cP}{HTML}{CC79A7}             % reddish purple P (P1 and the staged P)
\definecolor{cRef}{HTML}{4D4D4D}           % reference grey (arrows, neutral outlines)
\definecolor{cO}{HTML}{D55E00}             % vermillion: the flagged vacancy O site
\colorlet{cNeutral}{black!45}

% ---- page geometry (docs/venue.md: MDPI \textwidth 138.6 mm) ----
\newlength{\figwidth}\setlength{\figwidth}{138.6mm}

% ---- styles ----
\tikzset{
  every picture/.style={line width=0.6pt, font=\fontsize{8}{9.6}\selectfont},
  >={Stealth[length=2.2mm,width=1.6mm]},
  flow/.style={->, draw=cRef, line width=0.6pt},
  inject/.style={->, draw=cD, line width=0.9pt},
  box/.style={draw=cNeutral, rounded corners=1.2pt, minimum height=6mm, inner xsep=2.2pt,
              inner ysep=1.6pt, align=center, fill=white},
  model/.style={box, draw=#1, fill=#1!14, line width=0.8pt},
  modelname/.style={font=\fontsize{9}{10}\selectfont\bfseries, text=#1},
  quantity/.style={box, draw=cNeutral, fill=black!4},
  v2/.style={draw=#1, dashed, dash pattern=on 2pt off 1.2pt, line width=0.8pt, rounded corners=1.2pt},
  panel/.style={font=\fontsize{9}{10}\selectfont\bfseries, anchor=north west, inner sep=0pt},
  atom/.style={circle, inner sep=0pt, minimum size=#1},
}
\newcommand{\Ef}{$E_\mathrm{f}$}
 after
% \documentclass[tikz,border=1pt]{standalone}.
% Palette = src/dftgnn/viz/style.py COLORS (Okabe-Ito hues), extended with P and the v2 style.
\usepackage[T1]{fontenc}
\usepackage[scale=1.0]{tgheros}            % Helvetica-metric sans, Type 1, embeddable
\renewcommand{\familydefault}{\sfdefault}
\usepackage[eulergreek]{sansmath}
\sansmath
\usepackage{tikz}
\usetikzlibrary{positioning,arrows.meta,fit,backgrounds,calc,shapes.geometric,decorations.pathreplacing}

% ---- palette (hex as in style.py) ----
\definecolor{cB0}{HTML}{7F7F7F}            % grey   B0 physics floor
\definecolor{cRF}{HTML}{E69F00}            % orange RF-Kumagai
\definecolor{cS}{HTML}{0072B2}             % blue   S
\definecolor{cD}{HTML}{009E73}             % green  D
\definecolor{cP}{HTML}{CC79A7}             % reddish purple P (P1 and the staged P)
\definecolor{cRef}{HTML}{4D4D4D}           % reference grey (arrows, neutral outlines)
\definecolor{cO}{HTML}{D55E00}             % vermillion: the flagged vacancy O site
\colorlet{cNeutral}{black!45}

% ---- page geometry (docs/venue.md: MDPI \textwidth 138.6 mm) ----
\newlength{\figwidth}\setlength{\figwidth}{138.6mm}

% ---- styles ----
\tikzset{
  every picture/.style={line width=0.6pt, font=\fontsize{8}{9.6}\selectfont},
  >={Stealth[length=2.2mm,width=1.6mm]},
  flow/.style={->, draw=cRef, line width=0.6pt},
  inject/.style={->, draw=cD, line width=0.9pt},
  box/.style={draw=cNeutral, rounded corners=1.2pt, minimum height=6mm, inner xsep=2.2pt,
              inner ysep=1.6pt, align=center, fill=white},
  model/.style={box, draw=#1, fill=#1!14, line width=0.8pt},
  modelname/.style={font=\fontsize{9}{10}\selectfont\bfseries, text=#1},
  quantity/.style={box, draw=cNeutral, fill=black!4},
  v2/.style={draw=#1, dashed, dash pattern=on 2pt off 1.2pt, line width=0.8pt, rounded corners=1.2pt},
  panel/.style={font=\fontsize{9}{10}\selectfont\bfseries, anchor=north west, inner sep=0pt},
  atom/.style={circle, inner sep=0pt, minimum size=#1},
}
\newcommand{\Ef}{$E_\mathrm{f}$}


\begin{document}
\begin{tikzpicture}
  \tikzset{
    hs/.style={draw=black!60, line width=0.5pt, inner sep=0pt, minimum height=4.2mm, anchor=west},
    tst/.style={hs, fill=black!60, text=white},
    trn/.style={hs, fill=black!12},
    val/.style={hs, fill=black!30},
    title/.style={anchor=base west, inner sep=0pt},
    rowlab/.style={anchor=east, inner sep=2pt},
  }
  % two columns at x = 0 and 7.0; bars start 0.95 in and are 5.6 long
  \def\bx{0.95}

  % ---------------- (a) grouped hold-out, nested budgets ----------------
  \begin{scope}[shift={(0,0)}]
    \node[panel, anchor=base west] at (0,0.10) {(a)};
    \node[title] at (0.52,0.10) {grouped hold-out};
    \node[tst, minimum width=1.12cm] at (\bx,-0.35) {test};
    \node[trn, minimum width=4.48cm] at (\bx+1.12,-0.35) {pool};
    \foreach \b/\w [count=\k] in {25/1.15, 50/1.70, 100/2.30, 200/2.95, 400/3.70, 654/4.48} {
      \node[trn, minimum height=3.2mm, minimum width=\w cm] (b\k) at (\bx+1.12,-0.48-0.40*\k) {};
      \node[anchor=east, inner sep=2pt, font=\fontsize{7}{8}\selectfont] at (b\k.east) {$B = \b$};
    }
  \end{scope}

  % ---------------- (b) Kiyohara split ----------------
  \begin{scope}[shift={(7.0,0)}]
    \node[panel, anchor=base west] at (0,0.10) {(b)};
    \node[title] at (0.52,0.10) {Kiyohara split};
    \node[trn, minimum width=3.90cm] at (\bx,-0.35) {train};
    \node[val, minimum width=0.83cm] at (\bx+3.90,-0.35) {val};
    \node[tst, minimum width=0.87cm] at (\bx+4.73,-0.35) {test};
  \end{scope}

  % ---------------- (c) K-fold CV ----------------
  \begin{scope}[shift={(7.0,-1.15)}]
    \node[panel, anchor=base west] at (0,0.10) {(c)};
    \node[title] at (0.52,0.10) {$K$-fold CV};
    \def\sw{0.9333}
    \foreach \r/\lab [count=\k] in {1/fold 1, 2/fold 2, 3/fold 3, 6/fold $K$} {
      \pgfmathsetmacro{\y}{-0.35-0.50*(\k-1) - (\k==4 ? 0.48 : 0)}
      \node[rowlab] at (\bx,\y) {\lab};
      \foreach \s in {1,...,6} {
        \ifnum\s=\r \node[tst, minimum width=\sw cm] at ({\bx+(\s-1)*\sw},\y) {};
        \else \node[trn, minimum width=\sw cm] at ({\bx+(\s-1)*\sw},\y) {};\fi
      }
    }
    \node[inner sep=0pt] at (\bx-0.45,-1.76) {$\vdots$};
  \end{scope}

  % ---------------- (d) leave-one-chemistry-out ----------------
  \begin{scope}[shift={(0,-3.55)}]
    \node[panel, anchor=base west] at (0,0.10) {(d)};
    \node[title] at (0.52,0.10) {leave-one-chemistry-out};
    % chemistry-family blocks (widths in cm, sum 5.6) and the fold that holds each one out
    \def\fams{0.56/1, 0.27/2, 0.43/3, 0.20/4, 0.50/5, 0.32/2, 0.38/4, 0.24/1, 0.45/3, 0.28/5, 0.36/2, 0.25/4, 0.40/1, 0.22/3, 0.29/5, 0.45/4}
    \foreach \f [count=\k] in {1,...,5} {
      \pgfmathsetmacro{\y}{-0.35-0.50*(\k-1)}
      \node[rowlab] at (\bx,\y) {fold \f};
      \def\x{\bx}
      \foreach \w/\g in \fams {
        \ifnum\g=\f \node[tst, minimum width=\w cm] at (\x,\y) {};
        \else \node[trn, minimum width=\w cm] at (\x,\y) {};\fi
        \pgfmathsetmacro{\xn}{\x+\w}\global\let\x\xn
      }
    }
  \end{scope}
\end{tikzpicture}
\end{document}
